\scanpage{18}

% top, ms folio
\mspage{\textcircled{14}}

% l. 1, heading
\uline{Graphes cubiques \inserted{$C$} à sommet marqué} $u$

% ll. 2–6
Revient à la donnée d'un ens $\unsure{P|Q}_u$ de cardinal $5$ (l'ens des $5$ triangles de sommet $u\unsure{\}|)}$, ou encore des $5$ couples de pts mutuellt liés de $E(u)$), et d'un torseur $P$ sous $\Psi=\Psi(P_u,\mathbb F_2):\operatorname{Ker}\bigl(\mathbb F_2^{P_u}\xrightarrow[\inserted{(x_i)_{i\in I}\mapsto\sum x_i}]{}\mathbb F_2\bigr)$. \deleted{Ce torseur} \inserted{On \unsure{réconstitue|reconstruit|reconstitue}}

% ll. 7–9
$E(u)$ comme l'ens somme des $P\wedge_\Psi(\mathbb F_2,p_i)$, où les $p_i:\Psi\longrightarrow\mathbb F_2$ \inserted{$(i\in E)$} sont induits par les projections $\mathbb F_2^{P_u}\xrightarrow{\mathrm{pr}_i}\mathbb F_2$ $(i\in P_u)$, et $C-E(u)$ comme $P$, donc

% l. 10
\[
 C\simeq\Bigl(\coprod_{i\in E}P\wedge_\Psi(\mathbb F_2,p_i)\Bigr)\amalg P\amalg\{u\}.
\]

% ll. 10–20
La \unsure{recette|vieille} pour ce dictionnaire est la suivante: $E(u)$ est un ens de cardinal $10$ muni d'un autom involutif d'ordre $2$ ss pt fixe (car tt él de $E(u)$ est lié à un et un seul él. de $E(u)$, grâce à l'axiome du triangle), or une telle structure est déterminée par un ens $P_u$ de cardinal $5$ ($P_u=E(u)/\sigma$) et une famille d'ens à $2$ él. i.e. de $\mathbb F_2$ torseurs, paramétrés par $P_u$ -- i.e un torseur $Q$ sous $(\mathbb F_2)^{P_u}$.

% ll. 20–24
D'autre part, tt $x\in C-E(u)$ définit une partie de $E(u)$ -- savoir $E(x)\cap E(u)$ -- telle que \deleted{\unsure{sa}} intersection avec tte classe sous $\{e,\sigma\}$ i.e. tte fibre de $\overset{\inserted{=\widetilde Q_u}}{E(u)}\longrightarrow\unsure{Q|P}_u$ i.e. tte orbite $\subset E(u)$, soit de cardinal $1$.

% ll. 24–28
Or l'ens de ces parties n'est autre que $Q$ (qui est de cardinal $2^5=32$) alors que $C-E(u)$ est de card $16=32/2$. On trouve que $C-E(u)\longrightarrow Q$ est \uline{injectif}, et son image \inserted{$P$} est un torseur sous $\Psi\subset\mathbb F_2^{P_u}$, $\bigl(\Psi=\operatorname{Ker}(\mathbb F_2^{P_u}\longrightarrow\mathbb F_2)\bigr)$ -- donc on a sur $Q$ une reconstitution du groupe

% bottom left
% archive sheet number: 17
